\documentclass[11pt]{article}

\usepackage[T1]{fontenc}
\usepackage{lmodern}
\usepackage{microtype}
\usepackage[a4paper,margin=27mm,headheight=14pt]{geometry}
\usepackage{amsmath,amssymb,amsthm,mathtools}
\usepackage{booktabs,longtable,array}
\usepackage{enumitem}
\usepackage{xcolor}
\usepackage{hyperref}
\usepackage{fancyhdr}

\setlength{\emergencystretch}{2em}
\newcolumntype{L}[1]{>{\raggedright\arraybackslash}p{#1}}

\ifdefined\pdfinfoomitdate
  \pdfinfoomitdate=1
\fi
\ifdefined\pdftrailerid
  \pdftrailerid{}
\fi
\ifdefined\pdfsuppressptexinfo
  \pdfsuppressptexinfo=15
\fi

\hypersetup{
  colorlinks=true,
  linkcolor=blue!55!black,
  citecolor=blue!55!black,
  urlcolor=blue!55!black,
  pdftitle={Dependency-Colored Concentration for Finite-Alphabet Shared-Exclusions PID},
  pdfsubject={Finite-sample theorem, drift extension, local modulus, and formal evidence boundary},
  pdfkeywords={partial information decomposition, shared exclusions, dependency coloring, concentration},
  pdfauthor={pid-rs project analysis}
}

\pagestyle{fancy}
\fancyhf{}
\fancyhead[L]{pid-rs dependency-colored SxPID}
\fancyhead[R]{project validation}
\fancyfoot[C]{\thepage}

\newtheorem{theorem}{Theorem}[section]
\newtheorem{corollary}[theorem]{Corollary}
\newtheorem{lemma}[theorem]{Lemma}
\newtheorem{proposition}[theorem]{Proposition}
\theoremstyle{definition}
\newtheorem{definition}[theorem]{Definition}
\newtheorem{remark}[theorem]{Remark}

\newcommand{\cZ}{\mathcal{Z}}
\newcommand{\cC}{\mathcal{C}}
\newcommand{\Phat}{\widehat{P}}
\newcommand{\Pbar}{\overline{P}}
\newcommand{\Prb}{\mathbb{P}}
\newcommand{\E}{\mathbb{E}}
\newcommand{\Ind}{\mathbf{1}}
\newcommand{\supp}{\operatorname{supp}}
\newcommand{\normone}[1]{\left\lVert #1\right\rVert_1}
\newcommand{\sx}{\mathrm{sx}}

\title{\textbf{Dependency-Colored Concentration\\for Finite-Alphabet Shared-Exclusions PID}\\
\large Finite-sample law control, drift separation, local atom moduli,\\
and a formal evidence boundary}
\author{pid-rs project analysis}
\date{23 July 2026}

\begin{document}
\maketitle

\begin{abstract}
This document gives new project-defined validation for the Wibral-lineage categorical
shared-exclusions partial information decomposition (SxPID) in \texttt{pid-rs}.  The
paper-defined functional is not new.  The document derives a finite-sample concentration bound
for the complete empirical joint law when all complete rows share one common finite law and a
fixed coloring partitions observations into classes whose complete rows are mutually independent.
Dependence across classes can be arbitrary.  A
class-size proxy that is optimal within the declared H{\"o}lder--Hoeffding proof scheme improves
the ordinary color-count factor for unbalanced classes.  A telescoping error allocation
gives a time-uniform envelope.  A nonidentical-law extension separates random error from explicit
drift.  A stronger common-support modulus then transfers joint-law error to pointwise and averaged
SxPID atoms.  Lean checks the deterministic core, and an executable challenge suite retains
counterexamples to weaker assumptions.  The result does not validate continuous shared
exclusions, arbitrary time series, adaptive evaluation transforms, or binary64 asymptotics.
\end{abstract}

\begingroup
\small
\setlength{\parskip}{0pt}
\tableofcontents
\endgroup

\section{Status, scope, and provenance}

The mathematical functional, the concentration ingredients, and this repository artifact have
different origins.

\begin{center}
\begin{tabular}{L{0.25\textwidth}L{0.27\textwidth}L{0.38\textwidth}}
\toprule
Object & Origin & Status in this document \\
\midrule
Categorical SxPID & Paper-defined by Makkeh, Gutknecht, and Wibral; its lattice has a formal-logic
basis due to Gutknecht, Wibral, and Makkeh~\cite{makkeh2021,gutknecht2021} & Implemented for two to
four sources; receives a finite-sample bound and an almost-sure exact-real consistency implication
when the displayed envelope vanishes \\
Hoeffding bounded-variable inequality & Paper-defined classical result~\cite{hoeffding1963} & Used
inside each independent color class \\
Dependency-color concentration & Published background includes Janson's coloring method and later
H{\"o}lder formulations~\cite{janson2004,pelekis2017} & Project-defined proof composition with a
class-size proxy that is optimal within the declared H{\"o}lder--Hoeffding proof scheme; no novelty
claim \\
Empirical-law $L^1$ deviation & Published independent and identically distributed (i.i.d.)\
background includes Weissman
et al.~\cite{weissman2003} & One-sided subset union bound under the stated color contract \\
SxPID local modulus & Project-defined validation of the paper-defined functional & Safe exact-real
common-support baselines with exact M{\"o}bius row norms; no sharpness claim \\
Drift envelope & Project-defined extension & Concentration about the average row law plus explicit
bias to a reference law \\
\bottomrule
\end{tabular}
\end{center}

The machine-readable authority is \texttt{method-catalog.json}.  The rendered authority is
\texttt{METHODS.md}.  The Rust implementation routes covered by the final composition are the
stable categorical functions
\begin{center}
\nolinkurl{discrete_sxpid2}, \nolinkurl{discrete_sxpid2_averaged},
\nolinkurl{discrete_sxpid3}, \nolinkurl{discrete_sxpid3_averaged},
\nolinkurl{discrete_sxpid_n}, and \nolinkurl{discrete_sxpid_n_averaged}.
\end{center}
They are exported from \nolinkurl{stable::categorical}.  This document adds no estimator and no
public Rust API.  Its proof and executable artifacts are validation evidence.

\section{Finite setting and the dependence contract}

Let $Z_1,\ldots,Z_n$ take values in a finite alphabet $\cZ$ with $|\cZ|=K$.  For SxPID, one row is
the complete joint symbol
\begin{equation}
  Z_i=(S_{1i},\ldots,S_{mi},T_i).
\end{equation}
Dependence between source and target coordinates inside a row is unrestricted.  It is the signal
that PID analyzes.

For the common-law theorem, assume every row has marginal law $P$.  Let a deterministic color map
partition the indices into occupied classes $C_1,\ldots,C_r$.  Define
\begin{equation}
  n_j=|C_j|>0,
  \qquad
  n=\sum_{j=1}^{r}n_j.
\end{equation}
The complete rows $\{Z_i:i\in C_j\}$ must be mutually independent for every $j$.
Dependence across different classes can be arbitrary.  The color map must not depend on the
observed row values.

Define
\begin{equation}
  \Phat_n(z)=\frac1n\sum_{i=1}^{n}\Ind\{Z_i=z\},
  \qquad
  V_n=\left(\sum_{j=1}^{r}\sqrt{n_j}\right)^2,
  \qquad
  d_{\mathrm{eff},n}=\frac{V_n}{n}.
\end{equation}

\begin{proposition}[Effective-color range]\label{prop:effective-range}
The deterministic proxy obeys
\begin{equation}
  n\le V_n\le rn,
  \qquad
  1\le d_{\mathrm{eff},n}\le r.
\end{equation}
\end{proposition}

\begin{proof}
For the lower bound, expand the square of a sum of nonnegative square roots and retain its diagonal
terms.  For the upper bound, apply Cauchy--Schwarz:
\begin{equation}
  \left(\sum_j\sqrt{n_j}\right)^2
  \le\left(\sum_j1^2\right)\left(\sum_jn_j\right)=rn.
\end{equation}
The Lean artifact proves both inequalities over exact real numbers.
\end{proof}

The value $d_{\mathrm{eff},n}$ is a proof constant derived from the declared class sizes.  It is not an
estimated effective sample size.

\section{Finite-sample law concentration}

\begin{theorem}[Dependency-colored empirical-law bound]\label{thm:finite}
Assume $K\ge2$ and the contract in Section~2.  Then, for every $\varepsilon>0$,
\begin{equation}\label{eq:finite}
  \Prb\!\left(\normone{\Phat_n-P}\ge\varepsilon\right)
  \le
  \min\!\left\{
    1,
    (2^K-2)\exp\!\left(-\frac{n^2\varepsilon^2}{2V_n}\right)
  \right\}.
\end{equation}
If at most $d$ colors are occupied, then
\begin{equation}\label{eq:coarse}
  \Prb\!\left(\normone{\Phat_n-P}\ge\varepsilon\right)
  \le
  \min\!\left\{
    1,
    (2^K-2)\exp\!\left(-\frac{n\varepsilon^2}{2d}\right)
  \right\}.
\end{equation}
\end{theorem}

\begin{proof}
Fix a nonempty proper subset $A\subset\cZ$.  Put
\begin{equation}
  Y_i=\Ind\{Z_i\in A\}-P(A),
  \qquad
  S_j=\sum_{i\in C_j}Y_i.
\end{equation}
Let $R=\sum_j\sqrt{n_j}$ and select the H{\"o}lder exponents
\begin{equation}
  p_j=\frac{R}{\sqrt{n_j}},
  \qquad
  \sum_j\frac1{p_j}=1.
\end{equation}
Each $p_j\ge1$ because $R\ge\sqrt{n_j}$, so these are valid H{\"o}lder exponents.
For $\lambda>0$, generalized H{\"o}lder gives
\begin{align}
  \E\exp\!\left(\lambda\sum_jS_j\right)
  &\le
  \prod_j\left(\E e^{p_j\lambda S_j}\right)^{1/p_j}.
\end{align}
Rows are mutually independent inside each color.  Hoeffding's lemma for a centered random
variable with range length one gives
\begin{equation}
  \E e^{p_j\lambda S_j}
  \le
  \exp\!\left(\frac{n_jp_j^2\lambda^2}{8}\right).
\end{equation}
Therefore,
\begin{equation}\label{eq:mgf}
  \E\exp\!\left(\lambda\sum_jS_j\right)
  \le
  \exp\!\left(\frac{\lambda^2V_n}{8}\right).
\end{equation}
Chernoff's method, optimized at $\lambda=4nt/V_n$, yields
\begin{equation}\label{eq:subset-tail}
  \Prb\!\left(\Phat_n(A)-P(A)\ge t\right)
  \le
  \exp\!\left(-\frac{2n^2t^2}{V_n}\right).
\end{equation}
For any signed finite measure with total mass zero,
\begin{equation}\label{eq:tv}
  \frac12\normone{\Phat_n-P}
  =\max_{A\subseteq\cZ}\{\Phat_n(A)-P(A)\}.
\end{equation}
Set $t=\varepsilon/2$ and take a union bound over the $2^K-2$ nonempty proper subsets.  This proves
Equation~\eqref{eq:finite}.  Complement subsets already encode the opposite sign, so no additional
factor of two is required.  Proposition~\ref{prop:effective-range} gives
Equation~\eqref{eq:coarse}.
\end{proof}

For $K=1$, the $L^1$ distance is identically zero, and the finite-sample formulas reduce to the
true but empty statement $0\le0$ for every $\varepsilon>0$.  The time-uniform radius contains
$\log(2^K-2)$, so Section~4 requires $K\ge2$.
The subset factor is exponential in $K$, so the bound can be vacuous on a large alphabet.  This
result makes no sharpness or sample-complexity claim.  If $s=|\supp(P)|\ge2$ is known before the
data are inspected, then $2^s-2$ can replace $2^K-2$.  An observed support size is not a known
population support size.  When $s=1$, the law distance is zero and the result is vacuous.

\begin{proposition}[H{\"o}lder proxy optimum]\label{prop:holder-optimum}
For $q_j>0$ with $\sum_jq_j^{-1}=1$,
\begin{equation}
  V_n\le\sum_jn_jq_j.
\end{equation}
Equality holds for $q_j=R/\sqrt{n_j}$.
\end{proposition}

\begin{proof}
The exponent conditions force $q_j\ge1$ for every $j$.  Thus, they are valid H{\"o}lder
exponents.
Cauchy--Schwarz gives
\begin{equation}
  \left(\sum_j\sqrt{n_j}\right)^2
  =\left(\sum_j\sqrt{n_jq_j}\sqrt{q_j^{-1}}\right)^2
  \le\left(\sum_jn_jq_j\right)\left(\sum_jq_j^{-1}\right).
\end{equation}
The selected exponents make the two Cauchy--Schwarz vectors proportional.
\end{proof}

This is an optimum inside the H{\"o}lder--Hoeffding proof scheme for this partition.  It is not a
claim that $V_n$ is the best possible constant for every admissible joint law.

\section{Time-uniform envelope and consistency}

Assume $K\ge2$.  For an infinite sequence, require one fixed map from the positive integers to a finite or countable
color set.  Only finitely many colors can be occupied in each prefix.  The entire infinite
collection in each color must consist of mutually independent complete rows.  Let $V_n$ use the nonempty class
sizes in the first $n$ rows.  A uniform bound of at most $d$ occupied colors is required only for
the coarse radius in Equation~\eqref{eq:anytime-coarse}.

Fix $0<\alpha<1$ and define
\begin{equation}
  C_K=2^K-2,
  \qquad
  \alpha_n=\frac{\alpha}{n(n+1)}.
\end{equation}
The allocation has the exact finite identity
\begin{equation}\label{eq:telescope}
  \sum_{n=1}^{N}\alpha_n
  =\alpha\left(1-\frac1{N+1}\right)
  \le\alpha.
\end{equation}
Lean proves the unit-scale finite identity by induction.  Multiplication by the fixed scalar
$\alpha$ gives Equation~\eqref{eq:telescope}.

Define
\begin{align}
  R_n
  &=\sqrt{
      \frac{2V_n}{n^2}
      \log\!\left(\frac{C_Kn(n+1)}{\alpha}\right)
    }, \\
  D_n&=\min\{2,R_n\}.
\end{align}

\begin{theorem}[Anytime law envelope]\label{thm:anytime}
Under the infinite-sequence contract,
\begin{equation}\label{eq:anytime}
  \Prb\!\left(
    \forall n\ge1:\
    \normone{\Phat_n-P}\le D_n
  \right)\ge1-\alpha.
\end{equation}
\end{theorem}

\begin{proof}
For every $n$, substitution in Theorem~\ref{thm:finite} gives
\begin{equation}
  C_K\exp\!\left(-\frac{n^2R_n^2}{2V_n}\right)=\alpha_n.
\end{equation}
The Lean artifact checks the algebraic exponent cancellation.  Positivity, the square-root
substitution, and its probability interpretation remain in this prose proof.  A union bound and
Equation~\eqref{eq:telescope} give simultaneous control by $R_n$.  Every probability-law $L^1$
distance is also at most two.  Thus it is at most $D_n=\min\{2,R_n\}$.
\end{proof}

The coarser $d$-color radius is
\begin{equation}\label{eq:anytime-coarse}
  D_n^{(d)}=
  \min\!\left\{
    2,
    \sqrt{
      \frac{2d}{n}
      \log\!\left(\frac{C_Kn(n+1)}{\alpha}\right)
    }
  \right\}.
\end{equation}
Under one uniform fixed bound $d$, this is $O(\sqrt{d\log(n)/n})$.  For each positive rational
error threshold, the tails in Equation~\eqref{eq:coarse} are summable.  Borel--Cantelli and a
countable intersection over those thresholds therefore give
\begin{equation}
  \normone{\Phat_n-P}\longrightarrow0
  \quad\text{almost surely}.
\end{equation}
If the number of colors grows, the displayed envelope converges to zero under the sufficient
condition
\begin{equation}
  \frac{V_n\log n}{n^2}\longrightarrow0.
\end{equation}
This condition is not necessary for empirical-law convergence.  For example, i.i.d.\ rows still
converge when a needlessly fine singleton coloring gives $V_n=n^2$.  The condition is sufficient
for this declared-color envelope.  A singleton coloring also permits complete dependence across
rows, so the theorem cannot infer convergence from that coloring alone.

\begin{corollary}[Fixed-width overlapping windows]\label{cor:fixed-window}
Let $E_t$ be i.i.d.\ innovations.  Fix $w\in\mathbb{N}$ before evaluation, and let
\begin{equation}
  Z_t=f(E_t,\ldots,E_{t+w-1})
\end{equation}
for one fixed measurable finite-output map $f$.  Coloring row $t$ by $t\bmod w$ satisfies the
dependence contract with at most $d=w$ colors.
\end{corollary}

\begin{proof}
Two rows in the same color use disjoint innovation blocks.  Every finite collection in that color
is therefore mutually independent.  The fixed map and the i.i.d.\ innovation law give the common
row law.
\end{proof}

This corollary covers a fixed categorical overlapping-window construction from i.i.d.\
innovations.  It does not cover circular wraparound windows, a width selected from evaluation
outcomes, a rolling-start sequence of repeated claims, or innovations with unspecified time
dependence.

\begin{definition}[Strong $\ell$-dependence]
Let $\ell\in\mathbb{N}_0$.  An ordered sequence is strongly $\ell$-dependent when
\begin{equation}
  \sigma(Z_i:i\le t)
  \quad\text{is independent of}\quad
  \sigma(Z_i:i\ge t+\ell+1)
\end{equation}
for every $t$.
\end{definition}
Coloring index $t$ by $t\bmod(\ell+1)$ separates consecutive members of each color by more than
$\ell$.  Repeatedly separate the last row from the sigma-field generated by all preceding rows.
This ordered induction proves mutual independence of all complete rows in each color.  Pairwise
lag independence alone does not prove this property.  The common-law theorem still requires one
common row law.  If row laws differ, the drift theorem and its additional premises apply.

\section{Nonidentical row laws and explicit drift}

Assume $K\ge2$.  Let row $i$ have law $P_i$, and keep the within-color
mutual-independence contract.  Define
\begin{equation}
  \Pbar_n=\frac1n\sum_{i=1}^{n}P_i.
\end{equation}
Center the subset indicator for row $i$ by $P_i(A)$.  The H{\"o}lder--Hoeffding proof is unchanged.

\begin{theorem}[Average-law and reference-law bounds]\label{thm:drift}
For every $\varepsilon>0$,
\begin{equation}\label{eq:average-law}
  \Prb\!\left(\normone{\Phat_n-\Pbar_n}\ge\varepsilon\right)
  \le
  \min\!\left\{
    1,
    C_K\exp\!\left(-\frac{n^2\varepsilon^2}{2V_n}\right)
  \right\}.
\end{equation}
For an explicit reference law $P_\star$, put
\begin{equation}
  b_n=\normone{\Pbar_n-P_\star}.
\end{equation}
Then
\begin{equation}\label{eq:reference-law}
  \Prb\!\left(
    \normone{\Phat_n-P_\star}\ge b_n+\varepsilon
  \right)
  \le
  \min\!\left\{
    1,
    C_K\exp\!\left(-\frac{n^2\varepsilon^2}{2V_n}\right)
  \right\}.
\end{equation}
\end{theorem}

\begin{proof}
Equation~\eqref{eq:average-law} follows from the row-specific centering.  The triangle inequality
gives
\begin{equation}
  \normone{\Phat_n-P_\star}
  \le\normone{\Phat_n-\Pbar_n}+b_n,
\end{equation}
which proves Equation~\eqref{eq:reference-law}.
\end{proof}

If $\normone{P_i-P_\star}\le\beta_i$, then convexity gives
\begin{equation}
  b_n\le\frac1n\sum_{i=1}^{n}\beta_i.
\end{equation}
Under one fixed infinite row-law sequence and coloring with the same premises, the all-prefix
reference-law envelope is
\begin{equation}
  \Prb\!\left(
    \forall n\ge1:\
    \normone{\Phat_n-P_\star}\le\min\{2,b_n+D_n\}
  \right)\ge1-\alpha.
\end{equation}
The bias term is necessary.  Concentration about $\Pbar_n$ does not identify a fixed scientific
estimand.

\section{A stronger local shared-exclusions modulus}

Let $p$ and $q$ be two laws on the complete joint alphabet.  Put
\begin{equation}
  \delta=\normone{q-p},
  \qquad
  p_{\min}=\min_{z\in\supp(p)}p(z),
  \qquad
  L=\log(1/p_{\min}).
\end{equation}
Assume
\begin{equation}\label{eq:support-premise}
  \supp(q)\subseteq\supp(p),
  \qquad
  0\le\delta<2p_{\min}.
\end{equation}
Then the supports are equal.  For every shared-exclusions event $E$ that contains the keyed
supported realization,
\begin{equation}
  p(E)\ge p_{\min},
  \qquad
  |q(E)-p(E)|\le\frac\delta2.
\end{equation}
The factor one half follows from the variational identity; Lean proves the corresponding
zero-total-mass event lemma.

All logarithms are natural, so all information values use nats.  Fix a supported realization
$z=(s_1,\ldots,s_m,t)$ and a nonempty antichain $\alpha$ of nonempty source subsets.  For a source
subset $a$, the event $\{S_a=s_a\}$ fixes all source coordinates in $a$.  Define
\begin{align}
  A_\alpha(s)
    &=\bigcup_{a\in\alpha}\{S_a=s_a\}, \\
  B_\alpha(s,t)
    &=A_\alpha(s)\cap\{T=t\}, \\
  C(t)
    &=\{T=t\}.
\end{align}
Each event contains $z$.  For a law $Q$ whose three event masses are positive, define
\begin{align}
  c_\alpha^+(z;Q)
    &=-\log Q(A_\alpha(s)), \\
  c_\alpha^-(z;Q)
    &=\log\frac{Q(C(t))}{Q(B_\alpha(s,t))}, \\
  c_\alpha^\sx(z;Q)
    &=c_\alpha^+(z;Q)-c_\alpha^-(z;Q).
\end{align}
Let $M$ be the fixed M{\"o}bius matrix for the selected redundancy lattice.  For
$u\in\{+,-,\sx\}$, define the pointwise and averaged atoms
\begin{align}
  \pi_\alpha^u(z;Q)
    &=\sum_\beta M_{\alpha\beta}c_\beta^u(z;Q), \\
  \overline{\pi}_\alpha^u(Q)
    &=\sum_{z\in\supp(Q)}Q(z)\pi_\alpha^u(z;Q).
\end{align}

Define
\begin{equation}\label{eq:lambda}
  \Lambda(\delta,p_{\min})
  =-\log\!\left(1-\frac{\delta}{2p_{\min}}\right).
\end{equation}

\begin{lemma}[One-log modulus]\label{lem:one-log}
Under Equation~\eqref{eq:support-premise},
\begin{equation}
  |\log q(E)-\log p(E)|\le\Lambda(\delta,p_{\min}).
\end{equation}
\end{lemma}

\begin{proof}
Put $x=\delta/(2p_{\min})<1$.  Event control gives
\begin{equation}
  1-x\le\frac{q(E)}{p(E)}\le1+x.
\end{equation}
Monotonicity of the logarithm gives the lower side.  For the upper side,
\begin{equation}
  \log(1+x)\le x\le-\log(1-x).
\end{equation}
The Lean artifact proves this argument over exact real numbers.
\end{proof}

For one lattice node, the cumulative informative term contains one event logarithm.  The
misinformative term contains two.  Their difference is the net term.  Their errors are therefore
at most
\begin{equation}
  \Lambda,
  \qquad
  2\Lambda,
  \qquad
  3\Lambda.
\end{equation}
Let $M$ be the fixed M{\"o}bius matrix, and define the exact absolute row sum
\begin{equation}
  r_\alpha=\sum_\beta|M_{\alpha\beta}|.
\end{equation}
Lean proves the generic fixed-row inequality.  Hence the pointwise informative, misinformative,
and net atom errors are at most
\begin{equation}\label{eq:pointwise-moduli}
  r_\alpha\Lambda,
  \qquad
  2r_\alpha\Lambda,
  \qquad
  3r_\alpha\Lambda.
\end{equation}

Each cumulative informative and misinformative value under $p$ lies in $[0,L]$.  Each cumulative
net value lies in $[-L,L]$.  A weight-change decomposition therefore gives the common-support
averaged bounds
\begin{equation}\label{eq:averaged-moduli}
\begin{split}
  \text{informative:}\quad &r_\alpha\left(\Lambda+\frac{L\delta}{2}\right),\\
  \text{misinformative:}\quad &r_\alpha\left(2\Lambda+\frac{L\delta}{2}\right),\\
  \text{net:}\quad &r_\alpha(3\Lambda+L\delta).
\end{split}
\end{equation}
Equations~\eqref{eq:pointwise-moduli}--\eqref{eq:averaged-moduli} are conservative generic
M{\"o}bius-row baselines.  They do not use the paper-defined nonnegativity of the informative and
misinformative SxPID atoms.  They are valid upper bounds, but this paper does not claim that they
are sharp SxPID constants.

The informative and misinformative weight terms use $L\delta/2$.  A half-factor for the net term
does not follow from the generic cumulative-value range alone.  The centered half-range tightening
is new in \texttt{pid-rs}.  It is project-defined validation, not a claim of scientific novelty or
publication priority.

For completeness, the weight-change decomposition for one component is
\begin{align}
  \left|
    \sum_zq(z)\pi_\alpha^u(z;q)
    -\sum_zp(z)\pi_\alpha^u(z;p)
  \right|
  &\le
  \sum_zq(z)
    \left|\pi_\alpha^u(z;q)-\pi_\alpha^u(z;p)\right| \\
  &\quad+
  \left|\sum_z(q(z)-p(z))\pi_\alpha^u(z;p)\right|.
\end{align}
For any finite real function $f$ and signed weights $h$ with $\sum_z h(z)=0$, centering $f$ at
the midpoint of its range gives
\begin{equation}
  \left|\sum_zh(z)f(z)\right|
  \le
  \frac{\sum_z|h(z)|}{2}
  \left(\max_z f(z)-\min_z f(z)\right).
\end{equation}
Here $h=q-p$ and $\sum_z|h(z)|=\delta$.  At $p$, each cumulative informative and
misinformative value is in $[0,L]$.  Its M{\"o}bius atom therefore has oscillation at most
$r_\alpha L$, which gives $r_\alpha L\delta/2$.  Each cumulative net value is in $[-L,L]$.
Its atom has oscillation at most $2r_\alpha L$, which gives $r_\alpha L\delta$.  The net
calculation does not add separate informative and misinformative weight terms.

If $\delta\le p_{\min}/2$, then $x\le1/4$.  Lean also proves
\begin{equation}\label{eq:linearized-modulus}
  \Lambda
  \le\frac{x}{1-x}
  \le\frac{4x}{3}
  =\frac{2\delta}{3p_{\min}}.
\end{equation}
This improves the looser valid constants in the earlier finite-alphabet analysis.  All statements
in this section use exact real arithmetic.

\section{Composition with SxPID}

For common-law sampling, empirical support is a subset of population support almost surely.  On
the event in Theorem~\ref{thm:anytime}, the strict condition
\begin{equation}
  D_n<2p_{\min}
\end{equation}
forces empirical and population support to agree.  The realized error satisfies $\delta\le D_n$,
and $\Lambda(\delta,p_{\min})$ is nondecreasing in $\delta$.  Therefore, use $D_n$ on the
right-hand side of Equations~\eqref{eq:lambda}--\eqref{eq:averaged-moduli}.  This gives
simultaneous exact-real envelopes for every supported shared-exclusions cumulative term and every
fixed M{\"o}bius atom.

A usable numerical envelope needs a justified lower bound on $p_{\min}$.  The observed minimum
frequency is not such a bound.  A positive population cell can be absent from a finite sample.

For Theorem~\ref{thm:drift}, also require every $P_i$ to be supported inside $\supp(P_\star)$.
Define
\begin{equation}
  p_{\min,\star}=\min_{z\in\supp(P_\star)}P_\star(z).
\end{equation}
If
\begin{equation}
  b_n+D_n<2p_{\min,\star},
\end{equation}
then the realized error is at most $b_n+D_n$.  Monotonicity gives the same composition with
$b_n+D_n$ on the right-hand side.

A frozen finite-output transform can be handled conditionally on an independent training
artifact.  Conditional on that artifact, the color map must be fixed.  Every row must have the
stated conditional law.  The complete rows in every color class must be mutually independent.  A
random conditional support floor gives a conditional result unless a deterministic lower bound is
available.

The displayed common-law envelope proves almost-sure exact-real finite-alphabet SxPID plug-in
consistency under the
sufficient condition $V_n\log(n)/n^2\to0$; fixed $d$ is sufficient.  The displayed drift
envelope proves almost-sure exact-real reference-law SxPID consistency when this sufficient
condition holds and $b_n\to0$.  These are not necessary conditions under a stronger sampling
theorem.  The result does not
validate continuous $I_\cap^\sx$, arbitrary overlapping windows, adaptive preprocessing on
evaluation rows, unknown mixing rates, or binary64 asymptotics.

\section{Counterexamples and invalidated routes}

\begin{longtable}{L{0.24\textwidth}L{0.43\textwidth}L{0.23\textwidth}}
\toprule
Invalid route & Counterexample & Required correction \\
\midrule
\endfirsthead
\toprule
Invalid route & Counterexample & Required correction \\
\midrule
\endhead
Pairwise independence inside one color & Let $U$ be uniform on $\mathbb F_2^4$.  Use the 15
nonzero linear forms $a\mathbin{\cdot}U$ as binary rows.  Every pair is independent.  All rows are
zero with probability $1/16$.  The empirical $L^1$ error is then one, while the false one-color
bound is $2e^{-15/2}\approx0.00111$. & Require mutual independence of all complete rows in each
color. \\
\addlinespace
Remove the color factor & Take $m$ independent fair bits and copy each bit once into each of $d$
colors.  Each color is i.i.d., but only $m=n/d$ independent values remain. & Keep $V_n$, or keep
the coarser factor $d$. \\
\addlinespace
Infer convergence from one color per row & Set every row equal to one common fair bit.  Singleton
colors satisfy their local premise, but $V_n=n^2$ and the empirical law does not converge to the
fair law. & This coloring cannot prove convergence; use a stronger coloring or another theorem. \\
\addlinespace
Read the envelope condition as necessary & Let the rows be i.i.d.\ but assign each row a separate
color.  Then $V_n=n^2$, although the empirical law converges by the strong law. & State
$V_n\log(n)/n^2\to0$ only as a sufficient condition for this envelope. \\
\addlinespace
Data-adaptive coloring without the complete conditional contract & Set every row equal to one fair
bit, then select one occupied color from that bit.  Conditional rows are degenerate, but their
conditional law is not the common unconditional law. & Require fixed colors, or prove the full
conditional law and independence premises. \\
\addlinespace
Substitute a mixing label for the color premise & Use a stationary fair two-state Markov chain
with flip probability $\eta=1/100$.  Its transition dependence decays geometrically.  For two
rows, the empirical $L^1$ error is one with probability $1-\eta=0.99$, above the false one-color
bound $2e^{-1}$. & Use a theorem with an explicit mixing coefficient and rate, or prove a valid
independence coloring. \\
\addlinespace
Pairwise lag independence as strong $\ell$-dependence & The finite-field construction has pairwise
independence and a joint failure. & Require the sigma-field definition.  Residues modulo
$\ell+1$ then give valid colors by ordered induction. \\
\addlinespace
Halve the net weight term from the generic range alone & On two points, let
$h=(-\delta/2,\delta/2)$ and a generic row value be
$f=(-r_\alpha L,r_\alpha L)$.  Then $|\sum hf|=r_\alpha L\delta$.  This row value is not shown to
be realizable by an SxPID atom when $r_\alpha>1$.  It proves no SxPID sharpness claim. & Keep the
generic range baseline unless a separate SxPID argument supplies a smaller range. \\
\addlinespace
Allow $\delta=2p_{\min}$ & Move all mass from a least-probable supported cell.  Its required logs
become undefined. & Keep $\delta<2p_{\min}$ strict. \\
\addlinespace
Control only separate marginals & Set $S_1=S_2=X$.  For $0<\rho<1$, at the supported realization
$(X,T)=(0,0)$, use fair binary $X,T$ marginals with diagonal cells $(1+\rho)/4$ and off-diagonal cells
$(1-\rho)/4$, then replace $\rho$ by $-\rho$.  Marginals stay fixed, but pointwise redundancy
changes by $\log((1+\rho)/(1-\rho))$. & Control the complete joint law. \\
\addlinespace
Allow new support with a linear average bound & Let $p(0,0,0)=1$, and let
$q(0,0,0)=1-\varepsilon$, $q(1,1,1)=\varepsilon$.  Then $\normone{q-p}=2\varepsilon$, while the
averaged redundancy is
$-(1-\varepsilon)\log(1-\varepsilon)-\varepsilon\log\varepsilon$. & Require common support or use
a nonlinear boundary modulus. \\
\addlinespace
Use observed minimum frequency as a population floor & An unobserved positive cell has empirical
frequency zero. & Use an external floor or a separate simultaneous lower-confidence result. \\
\bottomrule
\end{longtable}

The executable challenge suite enumerates the finite-field pairwise-independence construction,
copied-color construction, singleton-color construction, and adaptive-color construction.  It also
checks a support-deletion boundary, an unspecified-mixing construction, a generic net-weight range
extremizer, a univariate-marginal-only construction, and a new-support construction.  The remaining
checks cover the telescoping allocation, class-size inequalities, high-precision one-log cases, and
all displayed bounds on three committed two-source law pairs.

The fixed-window fixture uses independent fair innovations $U_t,V_t$ and
\begin{equation}
  S_{1,t}=U_t\mathbin{\mathrm{OR}}U_{t+1},
  \qquad
  S_{2,t}=V_t\mathbin{\mathrm{OR}}V_{t+1},
  \qquad
  T_t=S_{1,t}\mathbin{\mathrm{XOR}}S_{2,t}.
\end{equation}
Even and odd starts are the two colors.  Its exact joint count weights are $1:3:3:9$.  The
high-precision Decimal oracle supplies expected logarithmic values.  For stored logarithmic
constants and bounds, the Rust reconstruction tolerance is
$32\varepsilon_{\mathrm{mach}}\max(1,\lvert x_{\mathrm{Rust}}\rvert,
\lvert x_{\mathrm{oracle}}\rvert)$.  For categorical estimator outputs, the absolute comparison
ceiling is $32\varepsilon_{\mathrm{mach}}$ nats, where
$\varepsilon_{\mathrm{mach}}=2^{-52}$ for binary64.  The test also checks that row reversal gives
bit-identical output.

\section{Formal and executable evidence boundary}

The Lean project is pinned to one toolchain and one exact dependency manifest.  The repository
checker rejects admitted proof placeholders, builds the project, and replays its declarations with
Lean's kernel checker.

The new Lean module proves:
\begin{itemize}[leftmargin=*,nosep]
  \item the event-mass $\delta/2$ lemma and exact attainment by the positive-coordinate event;
  \item positivity under the strict support margin;
  \item the one-log modulus, its monotonicity, and its rational simplification;
  \item the effective-color proxy bounds and the normalized factor range;
  \item the generic absolute linear-row bound;
  \item generic finite weighted-average perturbation bounds, including the centered half-range
    form;
  \item the exact unit-scale finite telescoping allocation; and
  \item the algebraic radius-exponent cancellation.
\end{itemize}

Lean does not encode random variables or conditional laws.  It does not encode generalized
H{\"o}lder, Hoeffding's lemma, the Chernoff argument, or the probability union bound.  The drift
result and fixed-window independence argument also remain outside Lean.  Lean does not encode SxPID
events or the redundancy lattice.  The identification of its generic lemmas with SxPID atoms also
remains outside Lean.  Rust refinement and binary64 rounding remain outside Lean.  The probability
theorem is a prose proof.  Lean machine-checks only the deterministic obligations listed above.
The Python suite uses exact rational arithmetic for finite probability and count identities, plus
100-digit Decimal arithmetic for logarithms.  Transcendental logarithms are not exact rational
values.  From each local count-table pair, the Rust test independently reconstructs the two
empirical laws, $\delta$, $p_{\min}$, $\Lambda$, $L$, and every stored node bound before it checks
the implementation outputs.  It uses the scale-aware reconstruction tolerance and the absolute
categorical-output ceiling declared above.  These checks are bounded executable evidence.  They
are not a general proof, a global binary64 bound, or external review.

\section{Ecosystem use boundary}

\begin{center}
\begin{tabular}{L{0.16\textwidth}L{0.33\textwidth}L{0.41\textwidth}}
\toprule
Consumer & Result supplied here & Still required \\
\midrule
Prisoma & Categorical SxPID law control after a frozen finite map when held-out rows meet the color
contract & Concrete row law, justified support floor, independent transform receipt, and a valid
window coloring \\
Galadriel & Dependence-aware categorical alternative to an i.i.d.\ plug-in claim & Removal of
same-window fitting and generic jitter routes; explicit observation and coloring contracts \\
Haldir & Mathematical target for block or residue-class designs & Direct dependency, corrected
atom interpretation, and proof of strong $\ell$-dependence or another valid coloring \\
Crebain & Finite-alphabet validation component for a future integration & Exact observation
mapping, run-log contract, direct dependency, and explicit scientific estimand \\
\bottomrule
\end{tabular}
\end{center}

This document does not assert that a consumer meets these premises.  Integration remains
\texttt{not\_claimed} until the consumer records and checks them.

\section{Reproduction}

From the repository root, run
\begin{verbatim}
python3 scripts/check-lean-finite-convergence.py
python3 scripts/generate-dependency-colored-sxpid-oracle.py
scripts/check-dependency-colored-sxpid-pdf.sh
\end{verbatim}
The PDF checker rebuilds this source with a fixed source-date epoch, rejects LaTeX diagnostics, and
compares the exact PDF bytes.  The Lean checker binds the toolchain, manifest, source set,
dependency revisions, dependency origins, and clean checkout state.  It removes ambient Git
routing and global or system configuration while retaining each dependency checkout's local
origin configuration for verification.

\section{Conclusion}

The result supplies a dependence-aware finite-sample path from a declared finite-row coloring to
exact-real categorical shared-exclusions PID bounds.  Its displayed envelope proves almost-sure
exact-real consistency under the sufficient condition $V_n\log(n)/n^2\to0$.  The corresponding
reference-law result also uses $b_n\to0$.  These are not necessary conditions under a stronger
sampling theorem.  The result separates
four objects that must remain separate: sampling fluctuation, deterministic drift, support margin,
and program arithmetic.  The color-size proxy improves the proof constant for unbalanced classes.
The telescoping allocation gives a simple time-uniform statement.  The local logarithmic modulus
gives a direct law-to-atom transfer once common support is justified.  Every weaker route listed
in Section~8 is retained as invalidated work.

\begin{thebibliography}{9}

\bibitem{makkeh2021}
A.~Makkeh, A.~J. Gutknecht, and M.~Wibral.
\newblock Introducing a differentiable measure of pointwise shared information.
\newblock \emph{Physical Review E}, 103:032149, 2021.
\newblock \url{https://doi.org/10.1103/PhysRevE.103.032149}.

\bibitem{gutknecht2021}
A.~J. Gutknecht, M.~Wibral, and A.~Makkeh.
\newblock Bits and pieces: understanding information decomposition from part-whole relationships
and formal logic.
\newblock \emph{Proceedings of the Royal Society A}, 477:20210110, 2021.
\newblock \url{https://doi.org/10.1098/rspa.2021.0110}.

\bibitem{hoeffding1963}
W.~Hoeffding.
\newblock Probability inequalities for sums of bounded random variables.
\newblock \emph{Journal of the American Statistical Association}, 58(301):13--30, 1963.
\newblock \url{https://doi.org/10.1080/01621459.1963.10500830}.

\bibitem{janson2004}
S.~Janson.
\newblock Large deviations for sums of partly dependent random variables.
\newblock \emph{Random Structures \& Algorithms}, 24(3):234--248, 2004.
\newblock \url{https://doi.org/10.1002/rsa.20008}.

\bibitem{pelekis2017}
C.~Pelekis, J.~Ramon, and Y.~Wang.
\newblock H{\"o}lder-type inequalities and their applications to concentration and correlation
bounds.
\newblock \emph{Indagationes Mathematicae}, 28(1):170--182, 2017.
\newblock \url{https://doi.org/10.1016/j.indag.2016.11.017}.

\bibitem{weissman2003}
T.~Weissman, E.~Ordentlich, G.~Seroussi, S.~Verd{\'u}, and M.~J. Weinberger.
\newblock Inequalities for the $L^1$ deviation of the empirical distribution.
\newblock HPL-2003-97 (R.1), 2003.
\newblock \url{https://shiftleft.com/mirrors/www.hpl.hp.com/techreports/2003/HPL-2003-97R1.pdf}.

\end{thebibliography}

\end{document}
